
\chapter*{Introduction générale}
\addcontentsline{toc}{chapter}{Introduction générale}

\par L'infrastructure informatique des grandes entreprises repose aujourd'hui sur trois piliers : le calcul, le stockage et le réseau. Pendant longtemps, chacun de ces piliers était géré par des équipements spécialisés distincts, ce qui entraînait une complexité administrative et des coûts élevés. L'hyperconvergence (HCI) est une architecture qui regroupe ces trois couches dans des serveurs standard, pilotés par une couche logicielle centralisée. Cette approche permet de simplifier l'administration, de réduire les coûts matériels et d'assurer la haute disponibilité de manière native. Dans un contexte de transformation numérique des entreprises industrielles, l'hyperconvergence est devenue une réponse concrète aux besoins d'infrastructures flexibles, évolutives et économiques.\\

\par Sonatrach, la compagnie nationale algérienne des hydrocarbures, dispose au sein de sa Direction Centrale Recherche \& Développement d'une infrastructure hyperconvergée reposant sur des solutions propriétaires, Nutanix et VMware vSAN, pour ses environnements de production. Si ces solutions offrent des garanties solides en termes de fiabilité et de support, elles présentent des contraintes structurelles importantes : des licences coûteuses dont le prix augmente chaque année, une dépendance forte vis-à-vis des éditeurs qui limite toute liberté d'adaptation, et un manque de souveraineté technique. Pour les environnements de test, de développement et de pré-production, qui nécessitent des clusters de trois à quatre n\oe{}uds seulement, déployer ces mêmes solutions reste disproportionné par rapport aux besoins réels. La question qui se pose est donc la suivante : comment concevoir et déployer une infrastructure hyperconvergée open source, capable de répondre aux exigences de fiabilité et de performance de Sonatrach, tout en constituant une alternative crédible et souveraine aux solutions propriétaires ?\\

\par Plusieurs solutions open source existent pour répondre à ce besoin. OpenStack est la plateforme cloud la plus complète et la plus standardisée, mais sa complexité de déploiement et sa consommation élevée de ressources la rendent inadaptée aux petits clusters. Harvester, développé par SUSE sur Kubernetes et KubeVirt, propose une interface moderne mais reste orienté vers les charges de travail conteneurisées et consomme des ressources importantes avant même d'allouer des capacités aux machines virtuelles classiques. Proxmox VE, en revanche, est un hyperviseur open source léger, qui déploie un cluster de trois n\oe{}uds en quelques heures et intègre nativement Ceph pour le stockage distribué. Pour la couche réseau, Open vSwitch avec des tunnels VLAN permet de créer une segmentation réseau complète sans nécessiter de modifications sur le switch physique. Ces technologies constituent aujourd'hui une pile technique mature et bien documentée, mais aucune d'elles n'intègre nativement un mécanisme d'équilibrage dynamique de charge des ressources du Clusters.\\

\par Dans le cadre de ce projet, nous avons conçu et déployé une infrastructure Cloud privée hyperconvergée entièrement open source sur trois serveurs Dell PowerEdge R720 fournis par Sonatrach. Nous avons installé Proxmox VE en mode bare metal sur chaque n\oe{}ud, puis formé un cluster haute disponibilité géré par Corosync, avec basculement automatique des machines virtuelles en cas de panne. Nous avons déployé un cluster Ceph à douze OSDs avec un facteur de réplication de trois, garantissant la persistance des données même en cas de défaillance d'un n\oe{}ud entier. Pour la couche réseau, nous avons conçu une architecture à neuf segments logiques, matérialisés par des tunnels VXLAN gérés par Open vSwitch, couvrant les flux d'infrastructure (Corosync, réplication Ceph, migration à chaud) et les zones applicatives (Dev, Prod, DMZ --- ``zone démilitarisée''). Nous avons également développé le \textbf{CLBS} (``Central Load Balancing Scheduler''), un service d'équilibrage dynamique de charge écrit en Go, qui surveille en continu la charge CPU et RAM des n\oe{}uds, calcule une bande de charge modérée et déclenche automatiquement des migrations à chaud pour redistribuer les machines virtuelles. Ce mécanisme est une adaptation de l'algorithme CSLB (``Central Scheduler Load Balancing''), étendu pour prendre en compte la RAM, avec des protections anti-oscillation et un circuit breaker. Pour compléter la plateforme, nous avons déployé une stack de supervision complète (InfluxDB, Prometheus, Grafana, Loki), des pipelines CI/CD via GitLab, une automatisation d'infrastructure par Ansible et Rundeck, une synchronisation temporelle par Chrony, et un pare-feu OPNsense assurant le routage et le filtrage inter-réseaux.\\

\par Ce mémoire est organisé en quatre chapitres :
\begin{itemize}
    \item \textbf{Chapitre I} : Contexte du projet, infrastructure Sonatrach, problématique et objectifs.
    \item \textbf{Chapitre II} : État de l'art, comparaison des solutions et justification des choix technologiques.
    \item \textbf{Chapitre III} : Conception de l'architecture, dimensionnement physique, plan d'adressage et logique CLBS.
    \item \textbf{Chapitre IV} : Déploiement, configuration via PXE et Ansible, et validation de la plateforme.
\end{itemize}
